% p3-08-phase-transition

\section{P3 §8. The Symbolic Phase Transition Conjecture}

8. The Symbolic Phase Transition Conjecture

         Figure (fig-p3-ch8-phase). Phase-transition-conjecture triptych --- Order parameter / Lucas
                                       dominance / Ratner reference.

  8.1 Motivation and Context
  The interplay between Trinity coarse-graining and the Pellis fine-scale structure
  suggests a threshold phenomenon: for approximation precisions below a critical value
  $\varepsilon$\textasciicircum{}, Trinity monomials suffice; above $\varepsilon$\textasciicircum{}, the Pellis series adds essential information.
  This threshold depends on the arithmetic nature of the target and is formally analogous
  to KAM theory's distinction between resonant and Diophantine frequencies.

  The analogy is the following. In the KAM theorem, the distinction between
  quasi-periodic (Diophantine) and resonant (rational) frequencies governs the
  persistence of invariant tori under perturbation. In the GOLDEN BRIDGE (Trinity-Pellis Bridge) framework, the
  distinction between "Trinity-dominant" and "Pellis-dominant" constants governs the
  efficiency of coarse-grained symbolic approximation. Constants near the Trinity lattice
  T are symbolic analogues of resonant frequencies; constants far from T require the full
  Pellis hierarchy, analogously to Diophantine constants requiring fine-scale
  continued-fraction expansions.

  8.2 Statement of the Conjecture

        Conjecture 8.1 (Symbolic Phase Transition; Conjecture). Let x > 0 and $\varepsilon$ > 0.
        Define the Trinity approximation exponent

  muT(x) = limsupL $\rightarrow$ $\in$fty (log |RL(x)|)/(ell(lambda\textasciicircum{}*L)) $\cdot$ (1)/(log $\varphi$-1).

  There exists a critical exponent mu\textasciicircum{}* = 1 such that:

  - Pellis-dominant phase (muT(x) > mu\textasciicircum{}*): Trinity monomials approximate                                 x
  sub-optimally per bit; Pellis corrections are necessary at all precisions.
  - Trinity-dominant phase (muT(x) < mu\textasciicircum{}*): x lies near a Trinity attractor,
  well-approximated by low-complexity monomials.

  - Critical line (muT(x) = mu\textasciicircum{}*): scale-free approximation behaviour; the Trinity and
  Pellis representations are equally efficient per bit.

  The conjecture asserts that this trichotomy is well-defined and that mu\textasciicircum{}* = 1 is the
  universal critical exponent, independent of the specific target x (except for arithmetic
  reasons).

  8.3 Lucas Ring Elements Are Trinity-Dominant

        Proposition 8.2. For any x $\in$ Z[$\varphi$] cap T, muT(x) = 0.

        Proof. By Proposition 5.9, RL(x) = 0 exactly when T\textasciicircum{}* = x $\in$ T. Hence |RL(x)| = 0
        and muT(x) = -$\in$fty (or, more precisely, muT(x) = 0 in the sense that the
        numerator in the defining ratio achieves -$\in$fty trivially). □

  This means the Lucas ring is the strongest possible attractor within the
  Trinity-dominant phase. The sanctioned seeds F17, …, F21, L7, L8 all lie in Z[$\varphi$] and are
  therefore maximally Trinity-dominant.

  8.4 Connection to the Margulis--Dani Theorem and Ratner Theory
  If one replaces T by the orbit of a rational point under a unipotent flow on a
  homogeneous space, the distribution of |RL(x)| is governed by equidistribution
  theorems of Ratner (1991) and Shah (1996). Formalising the Trinity lattice as a
  homogeneous space --- i.e., identifying T with a quotient of a Lie group by a lattice --- is
  a natural direction for future work. Such a connection would promote the Symbolic
  Phase Transition Conjecture from a heuristic to a theorem.

  8.5 Computational Evidence (Preliminary)
  Section 9.3 describes the protocol for estimating muT(x) from finite data. Preliminary
  experiments suggest:

  - alpha QED-1 approx 137.036: near the critical line, muT approx 1.0 ± 0.1.
  - mupe approx 1836.15: Trinity-dominant at low complexity budgets, muT approx 0.7.
  - 2 approx 1.41421: Pellis-dominant (as expected, since 2 notin Q(5)), muT approx 1.3.

  These are illustrative, subject to the floating-point caveats of Section 11.3, and should
  not be read as confirmed results.

\subsection{References}

- Vasilev, D. \textit{Flos Aureus Monograph: Trinity Constants}. DOI \href{https://zenodo.org/doi/10.5281/zenodo.19227877}{10.5281/zenodo.19227877}.
- CODATA 2018 recommended values: NIST \href{https://physics.nist.gov/cuu/Constants/}{https://physics.nist.gov/cuu/Constants/}.
- Heyrovska, R. Golden-angle FSC publications --- Heyrovský Institute of Physical Chemistry.
